% =====================================================================
%  FIELD DYNAMICS SECTION  --  insert into main.tex
%  Suggested placement: immediately after "Goal 1: Gravity = Time = Eddies"
%  (which defines the eddy field and the force law), or as a standalone
%  section before "Testable Predictions Summary".
%  NOTE: requires one bibliography entry for the GW-speed bound, cited below
%  as \cite{LIGO2017GW170817} (GW170817 / GRB 170817A, |c_gw-c|/c < ~1e-15).
% =====================================================================

\section{Field Dynamics: The Eddy Wave Equation}
\label{sec:dynamics}

Sections~\ref{sec:dynamics} onward supply a component absent from the framework
as stated so far. The eddy field $\omega$ and the response of matter to it are
specified---the latter through the force law
$a = -c^2\nabla(\omega/H_0)^2/2$---but the \emph{self-dynamics} of the field, i.e.
how $\omega$ evolves in time, are not. Without them the field cannot support
propagating disturbances, so quantities that depend on wave propagation (notably
the LISA delay of Sec.~\ref{sec:predictions}) rest on a kinematic assumption rather
than a dynamical result. This section provides the minimal dynamical completion,
addresses the ``complete kinetic theory of the eddy background'' identified as
future work in the abstract, and shows that eddy ripples propagate at $c$ under a
single, explicitly stated assumption.

\subsection{The field and the missing dynamics}

Define the dimensionless eddy field
\begin{equation}
\varphi \equiv \frac{\omega}{H_0}, \qquad
\varphi_0(x) = x^{\,n(x)} \ \ \text{(equilibrium profile, Eq.~\ref{eq:n}--\ref{eq:obs})}.
\end{equation}
The force law $a = -\tfrac{c^2}{2}\nabla\varphi^2$ is a \emph{static} relation: given
the configuration $\varphi(\mathbf{x})$ it returns the acceleration of test matter,
but it does not determine the time evolution of $\varphi$. A propagating ripple
requires $\varphi$ to be able to depart from $\varphi_0$ and relax back---that is, it
requires a term encoding inertia (resistance to change in time) alongside the
stiffness (resistance to spatial bending) already implied by the force law.

\subsection{Minimal quadratic action}

The most general local action quadratic in the field and its first derivatives is
\begin{equation}
S = \int \! dt\,d^3x
\left[\ \frac{\mu}{2}\left(\partial_t\varphi\right)^2
      - \frac{\kappa}{2}\left(\nabla\varphi\right)^2
      - V(\varphi)\ \right],
\label{eq:action}
\end{equation}
where $\mu$ is the inertia coefficient (dimensions of mass per unit length),
$\kappa$ is the stiffness coefficient (dimensions of tension, i.e.\ force), and
$V(\varphi)$ is a restoring potential with minimum at the equilibrium profile
$\varphi_0$. The Euler--Lagrange equation is
\begin{equation}
\mu\,\partial_t^2\varphi \;-\; \kappa\,\nabla^2\varphi \;+\; V'(\varphi) \;=\; 0 .
\label{eq:eom}
\end{equation}

\subsection{Ripple speed and the propagation fork}

Linearising about the equilibrium, $\varphi = \varphi_0 + \delta\varphi$ with
$\delta\varphi \propto \exp[i(\mathbf{k}\cdot\mathbf{x} - \omega_k t)]$, and using that
$\varphi_0$ satisfies the static balance $\kappa\nabla^2\varphi_0 = V'(\varphi_0)$,
Eq.~\eqref{eq:eom} gives the dispersion relation
\begin{equation}
\mu\,\omega_k^2 \;=\; \kappa\,k^2 \;+\; V''(\varphi_0).
\label{eq:disp_general}
\end{equation}
At large $k$ (or for $V''=0$) the propagation speed is
\begin{equation}
\frac{\omega_k}{k} \;\longrightarrow\; \sqrt{\frac{\kappa}{\mu}} \;\equiv\; c_\text{eddy},
\label{eq:speed}
\end{equation}
the classical wave speed $\sqrt{\text{tension}/\text{linear density}}$ of a stretched
medium. Crucially, the static results of the preceding sections fix \emph{neither}
$\mu$ \emph{nor} $\kappa$ individually, and therefore do not determine
$c_\text{eddy}$: at this level the eddy ripple speed is a free parameter, not a
derived quantity. We state this explicitly rather than absorb it silently.

\subsection{The relativistic closure}

The force law already contains the speed of light: $a = -\tfrac{c^2}{2}\nabla\varphi^2$.
The minimal completion consistent with this is that $\varphi$ is a Lorentz scalar, i.e.\
that inertia and stiffness are locked by

\begin{center}
\fbox{\ \ \parbox{0.85\textwidth}{\centering
\textbf{Assumption A (relativistic eddy sector).}\\[2pt]
$\displaystyle \mu = \kappa/c^2 \quad\Longrightarrow\quad c_\text{eddy} = c.$}\ \ }
\end{center}

\noindent
Equivalently, the action~\eqref{eq:action} takes the Lorentz-invariant form
\begin{equation}
S = \int\! dt\,d^3x\ \frac{\kappa}{2}
\left[\frac{1}{c^2}\left(\partial_t\varphi\right)^2 - \left(\nabla\varphi\right)^2\right]
- V(\varphi),
\qquad
\kappa\,\Box\,\varphi + V'(\varphi) = 0,
\label{eq:covariant}
\end{equation}
with $\Box \equiv c^{-2}\partial_t^2 - \nabla^2$. Assumption~A is a hypothesis, not a
theorem: it is the minimal choice compatible with the $c$ already present in the force
law, and it is the unique choice yielding causal (subluminal-bounded, luminal-limit)
signal propagation with no superluminal modes. We adopt it and mark each result that
depends on it.

\subsection{Consequences}

Under Assumption~A the ripple dispersion relation~\eqref{eq:disp_general} becomes
\begin{equation}
\omega_k^2 = c^2 k^2 + m^2, \qquad m^2 \equiv \frac{c^2}{\kappa}\,V''(\varphi_0).
\label{eq:disp_rel}
\end{equation}
Two consequences follow.
\begin{itemize}
\item \textbf{Luminal massless limit.} For $V''\to 0$, $\omega_k = c\,k$ exactly:
eddy ripples are non-dispersive and propagate at $c$. This is consistent with the
multimessenger bound $|c_\text{gw}-c|/c \lesssim 10^{-15}$ from GW170817/GRB\,170817A
\cite{LIGO2017GW170817}.
\item \textbf{Gravitational waves as a dynamical output.} The eddy ripples are
gravitational waves. Their luminal local speed, combined with the tilt profile
$\beta(x)$, produces a Shapiro-type propagation delay for signals traversing curved
regions, placing the LISA delay prediction of Sec.~\ref{sec:predictions} on a
dynamical footing rather than a kinematic postulate.
\end{itemize}

\subsection{Consistency constraint on the inertia}

If the inertia is identified with the eddy spin-energy density---the physically natural
choice for a rotational field, and consistent with $E=mc^2$---then $\mu \propto \varphi^2$
is field dependent. Equation~\eqref{eq:speed} would then give
$c_\text{eddy}^2 = \kappa/\mu \propto \varphi^{-2}$, a \emph{position-dependent} wave
speed (since $\varphi=\varphi_0(x)$), which is excluded by the observed universality of
$c_\text{gw}$. Universality therefore requires the stiffness to carry the same field
dependence, $\kappa \propto \varphi^2$, so that the ratio $\kappa/\mu = c^2$ remains
constant everywhere. This is precisely the content of Assumption~A generalised to a
field-dependent normalisation: a spin-energy inertia is admissible \emph{only} if
paired with a stiffness of matching dependence. The two ingredients are thus not
independent---the observational constraint links them.

\subsection{What is assumed versus derived}

\begin{center}
\begin{tabular}{ll}
\toprule
Status & Statement \\
\midrule
Derived (given A) & Eddy field obeys a wave equation, Eq.~\eqref{eq:covariant} \\
Derived (given A) & Ripples propagate at $c$; massless limit luminal, non-dispersive \\
Derived (given A) & GW propagation is a dynamical consequence, underpinning LISA delay \\
\midrule
Assumed & Lorentz invariance of the eddy sector (Assumption~A) \\
Motivation & The $c$ already in the force law; required for causal propagation \\
\midrule
Open & Shape of $V(\varphi)$ (hence any GW mass term / low-$k$ dispersion) \\
Open & Microscopic derivation of Assumption~A from the eddy background \\
Open (unchanged) & The $O(1)$ coefficient in $a_0 = H_0 c\,x_0/\sqrt{3}$; the \\
                 & $\sqrt{3}$ prefactor is independent of the dynamics introduced here \\
\bottomrule
\end{tabular}
\end{center}

\noindent
The dynamics of this section close the gap between the static force law and
wave propagation, at the cost of one explicit assumption (relativistic eddy sector).
They do not alter the derivation or the open status of the MOND coefficient $a_0$,
which depends on the equilibrium profile and the $\sqrt{3}$ projection, not on the
field's time evolution.
